\section{Torus bundles and compact parameter spaces}
\label{sec:torus-criterion}

The remaining task is to verify the $MU$-vanishing hypothesis from a
bundle embedding and then remove the polyhedral assumption on the
input.  We first construct the torus data used in
\Cref{not:torus-bundle}.  The Chern-class calculation takes place on a
tube, so that its vanishing controls the relative point class rather
than only its image in absolute cohomology.

\subsection{A framed torus detecting the point class}

Fix $p\ge1$ and set $t=2p+1$.  We describe once and for all the
geometric bundle data used in the criterion.

\begin{lemma}[Torus and line bundles]
\label{lem:framed-torus}
There is a smooth embedding
\[
 L\cong(\mathbb T^2)^p\ \subset\ (-c,c)^{2p+1}
\]
with a framed real normal line bundle.  There are complex line
bundles $\lambda_1,\ldots,\lambda_p$ on $L$, each pulled back from
its indicated two-torus factor, such that
\[
 H=\bigoplus_{i=1}^p\lambda_i\ \subset\ L\times\mathbb C^{2p},
 \qquad
 \left\langle\prod_{i=1}^p c_1(\lambda_i),[L]\right\rangle=1
\]
for a choice of orientation on $L$.
\end{lemma}

\begin{proof}
Begin with a circle embedded in $\mathbb R^2$ with its outward
normal framing.  If $\mathbb T^k$ has been embedded in
$\mathbb R^{k+1}$ with a framed normal line, its inclusion into
$\mathbb R^{k+2}$ has a framed normal two-plane bundle.  The boundary
of a sufficiently small tubular disk bundle is consequently a
smooth embedded hypersurface
$\mathbb T^k\times S^1$ in $\mathbb R^{k+2}$.  Its outward normal
is a framing.  Iteration produces $\mathbb T^{2p}$ in
$\mathbb R^{2p+1}$, and translation and rescaling put it in the
specified cube.

Choose a smooth degree-one map $\mathbb T^2\to S^2=\mathbb{CP}^1$;
one may collapse the complement of an oriented disk and then take
a smooth approximation.  The orthogonal complement of the
tautological line over $\mathbb{CP}^1$ is a line subbundle of
$\mathbb C^2$ with first Chern number one.  Pull it back to
$\mathbb T^2$, and then to each indicated factor of $L$.  Taking
their direct sum gives the stated embedding.  The evaluation follows
from the product orientation and the Chern number on each factor.
\end{proof}

Retain the slot bundle $H_n$ and the locus $\mathcal T(f,Z)$ of
\Cref{not:torus-bundle}.  The next lemma is valid for every real $a$;
the range $0<a<1/16$ is needed later for compression.  Its
complementary-bundle argument is related to the ordinary Chern-class
obstruction in \cite[proof of Theorem~3.3 in the author preprint]{HirshbergPhillips2022}.
We supply the additional
$MU$ square-zero calculation and supported Thom-class argument here.


\begin{lemma}[The cobordism criterion]
\label{lem:torus-mu}
Fix positive integers $p,n$ and a real number $a$.
Let $P$ be a finite polyhedron, let $f:P\to C(2p+1,n)$ be PL, and
let $Z=Z_{an}(2p+1,n)$.  If the slot bundle on $\mathcal T(f,Z)$
embeds into a trivial complex bundle of rank $0\le K<2pn$, then
\[
 \omega_{MU}(f,Z)=0.
\]
\end{lemma}

\begin{proof}
Put $M=(2p+1)n$, $r=pn$, $D=P\times Z$, and $d=d_f$.
The rank bound will force the pulled-back top $MU$ Chern class of $H_n$
to vanish.  We need this vanishing on the preimage of a neighborhood
of $L^n$, where multiplication by the normal Thom class will recover
the relative point test.

Fix the normal framing of $L^n$, of real rank $n$, obtained by
taking products in \Cref{lem:framed-torus}.  A tubular neighborhood
retraction extends $H_n$ and its orthogonal projection
$\Pi:L^n\to M_{2r}(\mathbb C)$ to a neighborhood of $L^n$.

We first extend the embedding to the pullback of a whole tube.
On the compact closed subset $\mathcal T(f,Z)\subset D$, compose
the assumed embedding with the source projection $\Pi\circ d$.
This represents it by a continuous $K$-by-$2r$ matrix.  Extend the
real and imaginary parts of its entries to $D$ by the Tietze
extension theorem \cite[Theorem~1.10.5]{Tao2010}.  On the original test locus its restriction to
the source fibers has a uniform positive lower bound on unit
vectors, by compactness.  This lower bound persists on the pullback
of a sufficiently thin tube: otherwise a sequence of unit source
vectors over points tending to the central locus would converge to
a vector contradicting injectivity there.  If the central locus is
empty, compactness instead gives a tube with empty pullback.
We may thus choose an open tube $U$ with compact closure in
$(-c,c)^M$ and retraction $\rho:U\to L^n$ such that on $d^{-1}(U)$
the pulled-back bundle $d^*\rho^*H_n$ embeds into $\mathbb C^K$.

We use the complex orientation of $MU$ and its Chern and Thom
classes, with the direct-sum Whitney formula
\cite[Lecture 5, Proposition 6 and the following discussion;
Lecture 6, Remark 7]{Lurie2010}.  List the $r$ line summands of
$H_n$ and write
\[
 x_i=c_1^{MU}(\lambda_i)\in MU^2(L^n),\qquad 1\le i\le r,
 \qquad \alpha=\prod_{i=1}^r x_i.
\]
Here the indices include both the slot and the two-torus factor.
Each $x_i$ is pulled back from a two-dimensional torus.  Since $MU$
is connective, $MU^4(\mathbb T^2)=0$: filtering that torus by its
cells reduces this assertion to $MU^q(\mathrm{pt})=0$ for $q>0$.
Consequently $x_i^2=0$.

Pull these classes to $d^{-1}(U)$.  If this space is nonempty, the
image of the embedded rank-$r$ bundle has an orthogonal complement
$F$ in $d^{-1}(U)\times\mathbb C^K$.  Thus $F$ has rank $K-r<r$.
With a formal variable $v$ recording Chern index, the
Whitney formula and the square-zero identities give
\begin{equation}
 \label{eq:inverse-chern}
 c_v^{MU}(F)
    =\prod_{i=1}^r(1+d^*\rho^*x_i\,v)^{-1}
    =\prod_{i=1}^r(1-d^*\rho^*x_i\,v).
\end{equation}
The coefficient of $v^r$ on the left is zero by rank.  The
coefficient on the right is $(-1)^r d^*\rho^*\alpha$.  Hence
\begin{equation}
 \label{eq:euler-tube-zero}
 d^*\rho^*\alpha=0\quad\hbox{in }MU^{2r}(d^{-1}(U)).
\end{equation}
The same assertion is automatic if the tube pullback is empty.
This uses only the direct-sum Chern formula; no formula for tensor
products of lines is involved.

For the support pairs below, relative $MU$ means reduced cohomology
of homotopy cofibers.  The collared ambient tube pairs and the final
polyhedral pair also admit the literal quotient descriptions used here.
Choose a closed smaller tubular disk bundle $V$ with
$L^n\subset\operatorname{int}V$ and $V\subset U$.
The normal framing gives a Thom class
\[
 \tau\in MU^n(U,U\setminus\operatorname{int}V).
\]
Indeed the relative quotient is the Thom space of a trivial real
rank-$n$ bundle, whose Thom class is the $n$-fold suspension of the
unit.  This construction works for odd as well as even $n$
\cite[Lecture 5, Example 1]{Lurie2010}.  The relative product
\[
 \eta=(\rho^*\alpha)\tau
      \in MU^M(U,U\setminus\operatorname{int}V)
\]
extends by excision to a class supported in $V$ in
$MU^M(\mathbb R^M,\mathbb R^M\setminus\operatorname{int}V)$.
Its image in
\begin{equation}
 \label{eq:ambient-point-group}
 MU^M\bigl(\mathbb R^M,
                 \mathbb R^M\setminus(-c,c)^M\bigr)
       \cong\widetilde{MU}^M(S^M)
       \cong\pi_0MU=\mathbb Z
\end{equation}
is a generator, up to the sign of the normal orientation.
To check the integer, apply the ordinary orientation map
$MU\to H\mathbb Z$.  It induces an isomorphism on the group in
\Cref{eq:ambient-point-group}.  The ordinary image of $\eta$ has
evaluation
\[
 \left\langle\prod_{i=1}^r c_1(\lambda_i),[L^n]\right\rangle=1
\]
by \Cref{lem:framed-torus} and the product orientation.  Thus the
image of $\eta$ is exactly the $MU$ origin class up to sign.

We give the pullback with its relative sets to retain this supported
information.  On $d^{-1}(U)$ the pullback of $\eta$ vanishes, because
its absolute factor is zero by \Cref{eq:euler-tube-zero}.  Excision
for the smaller closed tube therefore gives zero in
\[
 MU^M\bigl(D,D\setminus d^{-1}(\operatorname{int}V)\bigr).
\]
This excision is valid after pullback: the closure of
$D\setminus d^{-1}(U)$ lies in the interior of
$D\setminus d^{-1}(\operatorname{int}V)$, since $V\subset U$ and
$V$ is compact.  In particular the excision isomorphisms commute
with $d$ and with the relative products used above.
As $V\subset(-c,c)^M$, the relative sets satisfy
\[
 A_c(f,Z)\subset D\setminus d^{-1}(\operatorname{int}V).
\]
The induced map on relative cohomology sends this zero class to
\begin{equation}
 \label{eq:point-test-relative-mu}
 MU^M\bigl(D,\{(x,z):\|f(x)-z\|_\infty\ge c\}\bigr).
\end{equation}
The image is, up to sign, the pullback of the origin class in
\Cref{eq:ambient-point-group}, so it is the $MU$ point test and is
zero.  Finally the pair in \Cref{eq:point-test-relative-mu} is a
finite polyhedral pair: $f$ is PL, $Z$ is a finite union of
coordinate cubes, and the relative condition is a finite union of
linear inequalities after subdivision.  Its relative cohomology
is therefore the reduced cohomology of the exact quotient in
\Cref{def:point-test}.  This proves the required vanishing.
\end{proof}

\subsection{Compact parameter spaces and exact factorization}

\begin{lemma}[Finite parameter reduction]
\label{lem:finite-parameter-reduction}
Fix positive integers $p,n$ and a real number $a$.
Let $Y$ be compact metrizable, let $f:Y\to C(2p+1,n)$ be
continuous, and put $Z=Z_{an}(2p+1,n)$.  Suppose the slot bundle on
$\mathcal T(f,Z)$ embeds into a trivial complex bundle of nonnegative integer rank $K$.
There are a finite polyhedron $P$, a continuous map $j:Y\to P$,
and a PL map $f_P:P\to C(2p+1,n)$ such that
\[
 f=f_P\circ j,
\]
and the slot bundle on $\mathcal T(f_P,Z)$ also embeds into a
trivial bundle of rank $K$.
\end{lemma}

\begin{proof}
If $Y$ is empty, take $P$ empty and the unique maps.  Otherwise put
$M=(2p+1)n$ and $r=pn$.
Represent the embedding on the compact test locus by a
$K$-by-$2r$ matrix $A$, by first projecting onto the fixed source
bundle.  Extend its entries to $Y\times Z$ by the Tietze theorem
\cite[Theorem~1.10.5]{Tao2010}.
Choose $\delta>0$ such that on the central test locus
\begin{equation}
 \label{eq:original-injectivity}
 \|A(y,z)v\|\ge4\delta\|v\|
 \quad\hbox{for }v\in(H_n)_{f(y)-z}.
\end{equation}
Compactness of its unit sphere bundle supplies this number.  If
the locus is empty the condition is vacuous and any $\delta>0$
may be used.

The real algebra on $Y\times Z$ generated by continuous real
functions on $Y$ and the Euclidean coordinates of $z$ contains the
constants and separates points.  By the real Stone--Weierstrass theorem
\cite[Theorem~1.10.18]{Tao2010},
the real and imaginary entries of $A$ can be uniformly approximated
by elements of this algebra.  Only finitely many functions on $Y$
occur in these approximants.  Rescale them into $[-1,1]$ and denote
them by $\varphi_1,\ldots,\varphi_q$.  Include every coordinate of
$f$ explicitly, by setting
\[
 j_0(y)=(f(y),\varphi_1(y),\ldots,\varphi_q(y))
       \in W=[-1,1]^{M+q}.
\]
The approximating matrix has the form $B(j_0(y),z)$, where $B$ is
a matrix of polynomials on $W\times Z$.  Choose the approximation
so that its operator-norm error is less than $\delta$ everywhere
on $Y\times Z$.  By \Cref{eq:original-injectivity},
\begin{equation}
 \label{eq:coordinate-injectivity}
 \|B(j_0(y),z)v\|\ge3\delta\|v\|
 \quad\hbox{on the original slot bundle}.
\end{equation}

Let $C=j_0(Y)$ and let $f_0:W\to[-1,1]^M$ be projection onto the
first $M$ coordinates.  The bundle over a point $(u,z)$ with
$f_0(u)-z\in L^n$ depends only on $f_0(u)-z$.  Consequently
\Cref{eq:coordinate-injectivity} gives the same bound at every
such point of $C\times Z$, regardless of the choice of a preimage
of $u$ in $Y$.

For completeness, consider the following compact bad set:
\[
 \begin{split}
 \mathcal B=\{(u,z,v):\;&u\in W,\ z\in Z,\ \|v\|=1,
               \ f_0(u)-z\in L^n,\\
              &\Pi(f_0(u)-z)v=v,
               \ \|B(u,z)v\|\le2\delta\}.
 \end{split}
\]
Here $\Pi$ is the continuous orthogonal projection onto $H_n$.
The projection of $\mathcal B$ onto $W$ is compact and disjoint
from $C$, by \Cref{eq:coordinate-injectivity}.  Its complement
is therefore a relative open neighborhood of $C$ in $W$.
Subdivide $W$ into sufficiently small closed cubes and take the
union $P$ of all cubes meeting $C$.  Finiteness and the positive
distance from $C$ to the bad projection ensure that $P$ is contained
in this neighborhood.  If the bad projection is empty, any such
subdivision suffices.  This construction is relative to $W$; it
does not extend the first $M$ coordinates beyond the unit cube.

Regard $j_0$ as a map $j:Y\to P$ and put $f_P=f_0|_P$.  The map
$f_P$ is coordinate-linear and hence PL.  On
$\mathcal T(f_P,Z)$, the restriction of $B(u,z)$ to the slot
bundle is bounded below by $2\delta$ on unit vectors, and thus
is the required continuous bundle embedding.  Finally
$f_P\circ j=f$ by the choice of the first $M$ coordinates.
\end{proof}

\subsection{Completion of the compact-input theorem}

The factorization just proved retains the cube map exactly, while
moving the auxiliary embedding data to a finite polyhedron.  We can
therefore apply boundary compression there and pull it back without
losing its pointwise preservation clause.


\begin{proof}[Proof of \Cref{thm:torus-compression}]
Apply \Cref{lem:finite-parameter-reduction} to obtain $j:Y\to P$
and $f_P$ with the same bundle embedding and $f=f_P\circ j$.
By \Cref{lem:torus-mu} the $MU$ point test of $f_P$ vanishes.
\Cref{thm:boundary-compression} gives the stated compression maps
on $P^u$ for arbitrarily large $u$.  Composing with $j^{\times u}$
gives the maps on $Y^u$.  The factorization of $f$ is exact, so the
boundary agreement also pulls back exactly.
\end{proof}

\begin{remark}[The strict rank threshold]
\label{rem:rank-threshold-sharp}
The bound $K<2pn$ in \Cref{thm:torus-compression} cannot in general
be replaced by $K\le2pn$.  Take $n=1$, let $Y=C(2p+1,1)$, and let
$f$ be the identity.  For every $0<a<1/16$ the sparse set is
$Z_a(2p+1,1)=\{0\}$, so the torus locus is $L$ and the canonical
inclusion $H\subset L\times\mathbb C^{2p}$ supplies an embedding
at the endpoint $K=2p$.

Suppose that, for some $u\ge1$, a continuous map
$G:B\to B$, where $B=C(2p+1,u)$, has at least one boundary slot
everywhere and fixes every original boundary slot.  Its image lies
in the total boundary $\partial B$.  For $x\in\partial B$, at least
one boundary slot of $x$ remains fixed throughout the straight-line
homotopy $(1-s)x+sG(x)$, $0\le s\le1$.  This homotopy therefore
stays in $\partial B$, showing that $G|_{\partial B}$ has degree
one.  But $G:B\to\partial B$ extends this restriction over the
cube, forcing its degree to be zero.  This contradiction rules out
such a map for every $u$.  In particular, the positive bound
$(a/8)u$ in \Cref{thm:torus-compression}, which requires at least
one boundary slot, is impossible at this endpoint.
\end{remark}
